\subsection{域的代数闭包}

定义3.——设K是一个域。所谓K的代数闭包，我们指的是K的任意一个既是代数的又是代数闭的扩张。

例子.—— * 1) 复数域C是实数域R的代数闭包（Gen.～Top., VIII, p.～100） *

\begin{enumerate}
\def\labelenumi{\arabic{enumi})}
\setcounter{enumi}{1}
\tightlist
\item
  设 K 为一个域，且 \(\Omega\) 为 K 的一个代数闭扩张。如果 \(\bar{K}\) 是 K 在 \(\Omega\) 中的相对代数闭包，则由V，第20页，命题2， \(\bar{K}\) 是 K 的一个代数闭包。* 特别地，全体代数数构成的域（V，第 20 页，习题 2）是有理数域 Q 的一个代数闭包。*
\end{enumerate}

命题6。——设 \(\Omega\) 是域 \(K\) 的一个扩张。要使 \(\Omega\) 是 \(K\) 的一个代数闭包，必要且充分的条件是：它应为代数的，并且 \(K[X]\) 中的每个非常数多项式应在 \(\Omega[X]\) 中分解为若干个一次因式的乘积。

该条件由 (AC) 知是必要的。反之，假设 \(\Omega\) 在K上是代数的，并且 K{[}X{]} 上的每一个非常数多项式在 \(\Omega[\mathbf{X}]\) 中是次数为 1 的因子的乘积。设 \(\Omega'\) 是 \(\Omega\) 的一个代数扩张，并且令 \(x \in \Omega'\) 。由于 \(x\) 在 \(\Omega\) 上是代数的，且 \(\Omega\) 在K上是代数的， \(x\) 在K上是代数的（V，第19页，命题3）。设 \(f\) 为 \(x\) 在 K 上的极小多项式。根据假设，多项式 \(f \in \mathbf{K}[\mathbf{X}]\) 在 \(\Omega[\mathbf{X}]\) 中分解为一次因子的乘积，因此 \(x \in \Omega\) 。于是我们有 \(\Omega' = \Omega\) ，且 \(\Omega\) 是代数闭的，因为它满足（ \(\mathrm{AC}'''\) ）。

注记 1. --- 如果 \(\Omega\) 在 \(K\) 上是代数的，并且 \(K[X]\) 的每个非常数多项式在 \(\Omega\) 中有一个根，则 \(\Omega\) 是 \(K\) 的一个代数闭包（V，第156页，习题20）。

命题7。——设 \(\Omega\) 是域 \(\mathbf{K}\) 的一个代数扩张。

\begin{enumerate}
\def\labelenumi{\alph{enumi})}
\item
  如果 \(\Omega\) 是代数闭的，那么 \(\mathbf{K}\) 的每个代数扩张都同构于 \(\Omega\) 的一个子扩张。
\item
  反之，假设 K 的每个有限次代数扩张都同构于 \(\Omega\) 的一个子扩张；则 \(\Omega\) 是代数闭的。
\end{enumerate}

断言a)由定理1（V，第20页）推出。现在假设b)的假设成立，并考虑一个非常数多项式 \(f \in K[X]\) 。设 E 为 f 的一个分裂域（V，第～21 页，命题 4）；由于 E 在 K 上是有限次代数扩张（V，第～18 页，定理 2），我们可以假设 E 是 \(\Omega\) 的一个子扩张。于是多项式 f 是 \(\Omega[X]\) 中若干个一次多项式的乘积，且命题 6 表明 \(\Omega\) 是代数闭的。

我们现在可以证明一个域代数闭包的存在性和唯一性（在同构意义下）。

定理 2（施泰尼茨）。——设 K 为一个域；则 K 存在一个代数闭包。若 \(\Omega\) 与 \(\Omega'\) 是 K 的两个代数闭包，则存在一个 \(\Omega\) 到 \(\Omega'\) 上的 K-同构。

由命题 6，K 的一个代数闭包无非是 K{[}X{]} 中所有非常数多项式之集的一个分裂扩张。因此定理 2 可由 V，第～21 页，命题 4 及 V，第～22 页，推论得出。

推论。——设 K 与 K' 为两个域，Ω 为 K 的一个代数闭包，Ω' 为 K' 的一个代数闭包。对于 K 到 K' 的每个同构 u，存在一个 Ω 到 Ω' 的同构 v 且 v 扩展 u。

只需将定理 2 应用于 \(\Omega\) 与 \((\Omega', u)\) 这两个 \(\mathbf{K}\) 的代数闭包。

注记。——2）在前一推论的记号下，一般存在 \(\Omega\) 的异于恒等映射的 K-自同构。因此，一般来说，同构 \(v\) （它将 \(\Omega\) 映到 \(\Omega'\) 上，且扩张从 K 到 K' 上的同构 \(u\) ）并不唯一。出于类似的原因，一般来说，对于同一族多项式 \((f_i)_{i \in I}\) ，一个分裂扩张 E 到另一个分裂扩张 E' 的同构不止一个。我们回顾，与此相反，对于完全闭包我们有唯一性（V，第～5页）。

\begin{enumerate}
\def\labelenumi{\arabic{enumi})}
\setcounter{enumi}{2}
\tightlist
\item
  设 K 为一个域，且 \(\Omega\) 为 K 的一个代数闭包。那么对于非常数多项式的一个族 \((f_{i})_{i\in I}\) （在 K{[}X{]} 中），可以给出其分裂扩张的如下构造：设 \(R_{i}\) 为多项式 \(f_{i}\) 在 \(\Omega\) 中的根集，并且令 \(R=\bigcup_{i\in I}R_{i}\) 。则 \(K(R)\) 是 \(\Omega\) 的唯一的子扩张，它是一个分裂
\end{enumerate}

扩张，用于 \((f_i)_{i\in I}\) （V，第～22页，命题5）。

\begin{enumerate}
\def\labelenumi{\arabic{enumi})}
\setcounter{enumi}{3}
\tightlist
\item
  设K是一个有限域，且 \(\Omega\) 是K的一个代数闭包。那么 \(\Omega\) 是无限的（V，第～20页，命题3）；由于K的每个有限次扩张都是有限域， \(\Omega\) 是K的一个无限次代数扩张。
\end{enumerate}

